% TEMPLATE-MILC-v3.tex - plantilla maestra UNICA de las guias MILC v3.
%
% La usan las cuatro familias de guias, cada una con su driver:
%   · Grados 6-11          -> scripts/build-guias-g11.py
%   · Bebras               -> scripts/build-guias-bebras.py
%   · Semillero            -> scripts/build-guias-semillero.py
%   · Territorio interior  -> scripts/build-guias-territorio-interior.py
%
% Antes habia tres copias casi identicas (~400 lineas iguales, 6 distintas):
% cada mejora habia que aplicarla tres veces, y en la practica no se hacia
% (el motor de recursos vivia solo en dos de ellas). Lo que varia entre
% familias esta parametrizado en 6 placeholders que cada driver llena
% (se nombran aqui SIN su sintaxis real: el builder reemplaza por texto plano,
% tambien dentro de los comentarios, y un valor multilinea rompe el preambulo):
%
%   HEADER_IZQ        encabezado izquierdo de pagina
%   HEADER_DER        encabezado derecho de pagina
%   PORTADA_ETIQUETA  rotulo sobre el numero grande (GRADO/MOMENTO/MODULO)
%   PORTADA_SERIE     linea de serie de la portada
%   PORTADA_ID        identificador de la guia en la portada
%   PORTADA_CIRCULO   el circulito de grado (°), o vacio si no aplica
%   PDF_TITULO        titulo en texto plano (sin \textbf ni comillas TeX) para
%                     los metadatos del PDF (/Title); lo produce lib_xelatex.texto_pdf
%
% Composicion (auditoria 2026-09-03): las bandas de estacion piden espacio
% (\Needspace*) y prohiben el salto justo despues (\nopagebreak), asi no quedan
% huerfanas al pie; las cajas largas (softbox, drawbox, infoband, puentebox) son
% `breakable`, asi una caja de media pagina ya no arrastra todo a la siguiente
% dejando la anterior al 40 %. Todo texto sobre mostaza o verde va en milcNegro
% (blanco sobre #E5B400 da 1,93:1 y sobre #58A923 2,95:1; ninguno pasa WCAG AA).
% El resto de claves las documenta content/guias/_SCHEMA.md. El script valida
% que no queden placeholders sin reemplazar antes de escribir el archivo final.

% Etiquetado PDF (latex-lab): /Lang, estructura y texto alternativo de las
% figuras, para lectores de pantalla. Debe ir ANTES de \documentclass.
\DocumentMetadata{lang=es-CO,tagging=on}
\documentclass[11pt,letterpaper]{article}
% headheight=24pt: la cabecera derecha lleva dos lineas (identidad de la guia
% y estacion actual). headsep se acorta en la misma medida para que el bloque
% de texto quede exactamente donde estaba con headheight=12pt.
\usepackage[margin=2.0cm,headheight=24pt,headsep=13pt]{geometry}
\usepackage[table]{xcolor}
\usepackage{fontspec}
\usepackage[spanish]{babel}
\usepackage{tikz}
% Librerías TikZ para los diagramas de las guías (bloque `recursos:`):
%  · babel  → babel en español vuelve activos caracteres como «>», y TikZ los usa
%             en sus opciones (>=stealth, ->). Sin esta librería, cualquier
%             diagrama con flechas revienta con «\language@active@arg> has an
%             extra }». Debe ir ANTES de que se dibuje nada.
%  · positioning        → sintaxis «below left=2pt and 3pt».
%  · decorations.pathreplacing → llaves ({brace}) para señalar rangos.
\usetikzlibrary{babel,positioning,decorations.pathreplacing}
\usepackage{graphicx}
% Circuitos: el trabajo de electrónica se hace hoy con micro:bit y sus sensores
% integrados (MakeCode, sin cableado), así que no se necesitan esquemáticos.
% Si algún día entran montajes con protoboard, basta descomentar esta línea:
% los diagramas .tex/.tikz declarados en `recursos:` ya se incrustan solos.
% \usepackage{circuitikz}
\usepackage[most]{tcolorbox}
\usepackage{tabularx}
\usepackage{array}
\usepackage{fancyhdr}
\usepackage{titlesec}
\usepackage[document]{ragged2e}  % bandera a la derecha en todo el cuerpo
\usepackage{enumitem}
\usepackage{microtype}
\usepackage{etoolbox}
\usepackage{needspace}
% hyperref al final del preambulo: metadatos del PDF (titulo, autor, idioma,
% familia), marcadores por estacion (\pdfbookmark) y enlaces sin recuadro.
\usepackage[hidelinks,
  pdftitle={Actuadores — luces y sonido con compás},
  pdfauthor={PhD. Álvaro Cárdenas Orozco},
  pdfsubject={Tecnología e Informática · Grado 8 · Guía 15 · MILC v3},
  pdflang={es-CO},
  bookmarksopen=true,bookmarksnumbered=false]{hyperref}

% Helvetica Neue no trae la flecha U+2192, asi que cada «→» escrito literalmente
% en el YAML se compilaba a NADA, en silencio (462 flechas en 61 guias: rutas del
% tipo «paleta → tipografia → lockup» salian rotas). Mapeamos el caracter a su
% version matematica, que si tiene glifo. Asi el autor puede seguir escribiendo
% «→» comodamente en el YAML.
\usepackage{newunicodechar}
\newunicodechar{→}{\ensuremath{\rightarrow}}
% Mismo problema con los simbolos de marca y vinetas: el «✓» de «una palomita ✓»
% salia como cuadrito vacio en el PDF de todas las guias que lo usaban.
\usepackage{pifont}
\newunicodechar{✓}{\ding{51}}
\newunicodechar{✗}{\ding{55}}
\newunicodechar{✔}{\ding{52}}
\newunicodechar{•}{\textbullet}
\newunicodechar{≥}{\ensuremath{\geq}}
\newunicodechar{≤}{\ensuremath{\leq}}

\setmainfont{Helvetica Neue}
\setsansfont{Helvetica Neue}
\setlength{\parindent}{0pt}
\setlength{\parskip}{6pt}
% Sin particion de palabras: en 6.º-7.º una palabra cortada por guion al final
% de linea («Comuni-cacion») frena la lectura mucho mas que una linea corta.
\hyphenpenalty=10000
\exhyphenpenalty=10000
\pretolerance=2000
\tolerance=3000
\setlength{\emergencystretch}{3em}
\linespread{1.10}
\renewcommand{\arraystretch}{1.25}
\setlist{leftmargin=5mm,itemsep=1mm,topsep=1mm}
\newcolumntype{Y}{>{\RaggedRight\arraybackslash}X}

\definecolor{milcMagenta}{HTML}{E6007E}
\definecolor{milcVino}{HTML}{5A0038}
\definecolor{milcTurquesa}{HTML}{0093A5}
\definecolor{milcVerde}{HTML}{58A923}
\definecolor{milcMostaza}{HTML}{E5B400}
\definecolor{milcOcre}{HTML}{B86600}
\definecolor{milcNegro}{HTML}{241816}
\definecolor{milcGris}{HTML}{F7F7F4}
\definecolor{milcLinea}{HTML}{E8E2DD}
\definecolor{milcGradePrimary}{HTML}{FF6600}
\definecolor{milcGradeDark}{HTML}{9A3E00}
\definecolor{milcGradeSoft}{HTML}{FFE4D0}
\definecolor{milcGradeAccent}{HTML}{CFE7FF}
\definecolor{milcGradeText}{HTML}{FFFFFF}
\color{milcNegro}

\pagestyle{fancy}
\fancyhf{}
% Cabecera de dos lineas: identidad de la guia arriba y, a la derecha y en
% gris, la estacion en curso (\markright desde \milcestacion). Antes de la
% primera estacion la segunda linea va vacia; el \strut mantiene la altura
% para que la linea de arriba no baile entre paginas.
\lhead{\small Tecnología e Informática · Grado 8\\[-1pt]{\footnotesize\strut}}
\rhead{\small Guía 15 · MILC v3\\[-1pt]{\footnotesize\strut\color{milcNegro!65}\rightmark}}
\cfoot{\small\thepage}
\renewcommand{\headrulewidth}{0pt}

% Sobre mostaza (#E5B400) y verde (#58A923) el texto va en milcNegro; sobre el
% resto de la paleta, en blanco. Se usa dentro de fonttitle/fontupper, donde un
% \color posterior manda sobre coltitle.
\newcommand{\milctintasobre}[1]{%
  \ifboolexpr{test{\ifstrequal{#1}{milcMostaza}} or test{\ifstrequal{#1}{milcVerde}}}%
    {\color{milcNegro}}{\color{white}}}

\newtcolorbox{titlebox}[1]{
  enhanced,colback=#1,colframe=#1,arc=7mm,boxrule=0pt,
  left=7mm,right=7mm,top=3mm,bottom=3mm,
  before skip=8pt,after skip=8pt,
  fontupper=\sffamily\bfseries\Large\color{white}
}

\newtcolorbox{softbox}[3]{
  enhanced,breakable,pad at break=3mm,
  colback=#2,colframe=#1,borderline west={4pt}{0pt}{#1},
  arc=2mm,boxrule=.4pt,left=4mm,right=4mm,top=3mm,bottom=3mm,
  before skip=6pt,after skip=8pt,title={#3},coltitle=white,
  colbacktitle=#1,fonttitle=\sffamily\bfseries\milctintasobre{#1},fontupper=\RaggedRight,
  attach boxed title to top left={xshift=3mm,yshift=-2mm},
  boxed title style={arc=2mm, boxrule=0pt}
}

\newtcolorbox{drawbox}[2]{
  enhanced,breakable,pad at break=4mm,
  colback=white,colframe=#1,arc=4mm,boxrule=1pt,
  left=5mm,right=5mm,top=4mm,bottom=4mm,before skip=8pt,after skip=8pt,
  title={#2},coltitle=white,colbacktitle=#1,fonttitle=\sffamily\bfseries\milctintasobre{#1},
  fontupper=\RaggedRight,
  attach boxed title to top left={xshift=4mm,yshift=-2mm},
  boxed title style={arc=3mm, boxrule=0pt}
}

\newtcolorbox{infoband}[1]{
  enhanced,breakable,pad at break=2mm,colback=#1!8,colframe=#1!50,arc=2mm,boxrule=.5pt,
  left=4mm,right=4mm,top=2mm,bottom=2mm,before skip=4pt,after skip=4pt,
  fontupper=\small\RaggedRight
}

% Puente narrativo entre secciones (contrato editorial Conéctate).
% Da continuidad: "Acabas de X. Ahora vas a Y porque Z."
% Visualmente: banda izquierda en color del grado, texto en itálica pequeña.
\newtcolorbox{puentebox}{
  enhanced,breakable,pad at break=2mm,colback=milcGradeSoft,colframe=milcGradeDark,
  borderline west={3pt}{0pt}{milcGradeDark},
  arc=0mm,boxrule=0pt,left=5mm,right=4mm,top=2.5mm,bottom=2.5mm,
  before skip=4pt,after skip=8pt,
  fontupper=\itshape\small\color{milcNegro}\RaggedRight
}
% La flecha va en modo matematico a proposito: Helvetica Neue NO tiene el glifo
% U+2192, asi que \textrightarrow se compilaba a NADA («Missing character: There
% is no → in font Helvetica Neue») y el puente salia sin flecha. En modo
% matematico la flecha viene de la fuente matematica y si se dibuja.
\newcommand{\puente}[1]{%
  \begin{puentebox}%
    \textcolor{milcGradeDark}{$\rightarrow$}\hspace{2mm}#1%
  \end{puentebox}%
}

\newcommand{\linead}[1]{\noindent\color{milcLinea}\rule{#1}{0.5pt}\color{milcNegro}}
\newcommand{\respuesta}[1]{\par\vspace{2mm}\foreach \n in {1,...,#1}{\noindent\color{milcLinea}\rule{\linewidth}{0.5pt}\par\vspace{5mm}}\color{milcNegro}}
\newcommand{\checkbox}{\textcolor{milcGradeDark}{\fbox{\phantom{x}}}\hspace{5pt}}

% ─── Listas aireadas para las guías socioemocionales ────────────────────
% Componer las verificaciones como «\checkbox texto\\» deja el texto que salta
% de línea debajo del cuadrito y apelmaza el bloque. Estos entornos alinean la
% continuación y airean los ítems. Los usa Territorio Interior; cualquier
% familia puede hacerlo.
\newlist{checklista}{itemize}{1}
\setlist[checklista]{
  label=\textcolor{milcGradeDark}{\fbox{\phantom{x}}},
  leftmargin=9mm, labelsep=3.2mm, itemsep=2.4mm,
  topsep=2.5mm, parsep=0pt, partopsep=0pt, before=\RaggedRight
}
\newlist{pasoslista}{enumerate}{1}
\setlist[pasoslista]{
  label=\textcolor{milcGradeDark}{\sffamily\bfseries\arabic*.},
  leftmargin=8mm, labelsep=3mm, itemsep=2.8mm,
  topsep=1mm, parsep=0pt, partopsep=0pt, before=\RaggedRight
}

\newcommand{\phasecard}[4]{
\begin{tcolorbox}[enhanced,colback=#3,colframe=#2,arc=3mm,boxrule=.5pt,height=3.05cm,valign=center,halign=center,left=1.5mm,right=1.5mm,top=2mm,bottom=2mm]
{\sffamily\bfseries\fontsize{22}{24}\selectfont\textcolor{#2}{#1}}\par
{\sffamily\bfseries\small #4}
\end{tcolorbox}
}

% ── Piezas del rediseno 2026 (legibilidad 6.º-11.º) ───────────────────
% Banda de estacion: reemplaza el titlebox macizo. Titulo a la izquierda,
% dato practico (tiempo/modalidad) a la derecha, en una sola linea.
% Nunca huerfana: pide 9 lineas libres (o salta de pagina rellenando con
% \vfill) y prohibe el salto justo despues de la banda. Nueve y no seis:
% la caja partible que sigue necesita ~80pt (titulo adosado, relleno y dos
% lineas) para arrancar en la misma pagina; con menos, tcolorbox la manda
% entera a la siguiente y la banda se queda sola. El penalty 10000 va
% en `after`, ANTES del espacio posterior: un `after skip` (aunque sea 0pt)
% es un \vskip tras la caja, y ese glue es un punto de corte legal que un
% \nopagebreak escrito despues de \end{tcolorbox} ya no cierra. Cada banda
% deja marcador PDF y actualiza la cabecera derecha.
\newcounter{milcestacion}
\newcommand{\milcestacion}[2]{%
\Needspace*{9\baselineskip}%
\stepcounter{milcestacion}%
\pdfbookmark[1]{#2}{milcestacion\themilcestacion}%
\markright{#2}%
\begin{tcolorbox}[enhanced,colback=#1,colframe=#1,arc=2mm,boxrule=0pt,
  left=5mm,right=5mm,top=2.6mm,bottom=2.6mm,before skip=11pt,
  after={\par\nopagebreak\vskip7pt}]
{\sffamily\bfseries\fontsize{15}{18}\selectfont\milctintasobre{#1}#2}
\end{tcolorbox}}

% Tarjeta de paso numerada: sustituye las tablas «Pilar / Aplicacion», que
% eran formato de consulta docente, no de estudiante. El numero va en un
% circulo sobre el borde para que la secuencia se lea de un vistazo.
\newcommand{\milcpaso}[3]{%
\Needspace{4\baselineskip}%
\begin{tcolorbox}[enhanced,colback=white,colframe=milcLinea,arc=1.5mm,boxrule=.9pt,
  left=14mm,right=4mm,top=3mm,bottom=3mm,before skip=4pt,after skip=4pt,
  fontupper=\RaggedRight,
  overlay={\node[anchor=north west,circle,fill=#2,text=white,inner sep=0pt,
    minimum size=7.4mm,font=\sffamily\bfseries\small]
    at ([xshift=3.6mm,yshift=-3.2mm]frame.north west) {#1};}]
#3
\end{tcolorbox}}

% Casilla marcable de tamano util con lapiz (la anterior era \fbox{\phantom{x}},
% de alto variable segun el texto que la seguia).
\newcommand{\milcmarca}{\framebox[3.8mm]{\rule{0pt}{3.4mm}}}

% Fila marcable de lista suelta (checks de la estacion 1).
\newcommand{\milccheckrow}[1]{%
\par\vspace{1.8mm}\noindent
\makebox[6.4mm][l]{\raisebox{-.55mm}{\textcolor{milcNegro}{\milcmarca}}}%
\parbox[t]{\dimexpr\linewidth-6.4mm\relax}{\RaggedRight #1}\par}

% Celda marcable dentro de la rubrica (tabla de autoevaluacion). Misma
% sangria francesa, pero pensada para vivir dentro de una columna X.
\newcommand{\milcopcion}[1]{%
\makebox[5.2mm][l]{\raisebox{-.55mm}{\textcolor{milcNegro}{\milcmarca}}}%
\parbox[t]{\dimexpr\linewidth-5.2mm\relax}{\RaggedRight #1}}

% ── Recursos visuales de la guía ──────────────────────────────────────
% Declarados en el bloque `recursos:` del YAML y emitidos por el builder.
% \guiaFigura   → imagen rasterizada o PDF (foto, carta celeste, montaje).
% \guiaDiagrama → fuente TikZ/circuitikz (.tex) incrustada como vector:
%                 es la vía para circuitos y esquemáticos editables.
% [#1] = texto alternativo (alt) · #2 = ruta relativa al .tex · #3 = pie
% El alt llega al PDF etiquetado (/Alt) y lo lee el lector de pantalla.
\newcommand{\guiaFigura}[3][]{%
\begin{center}
\includegraphics[width=0.88\linewidth,height=0.38\textheight,keepaspectratio,alt={#1}]{#2}\\[1.5mm]
{\footnotesize\itshape\textcolor{milcNegro!75}{#3}}
\end{center}
}

\newcommand{\guiaDiagrama}[2]{%
\begin{center}
\input{#1}\\[1.5mm]
{\footnotesize\itshape\textcolor{milcNegro!75}{#2}}
\end{center}
}

\begin{document}
\thispagestyle{empty}

% ── Portada ───────────────────────────────────────────────────────────
% Antes: un rectangulo de color a pagina completa. Se veia bien en pantalla,
% pero en la fotocopiadora del colegio se comia el toner y salia gris sucio.
% Ahora la banda ocupa el tercio superior y el resto va en blanco.
\begin{tikzpicture}[remember picture,overlay]
  \path[fill=milcGradePrimary]
    (current page.north west) rectangle ([yshift=-10.6cm]current page.north east);

  \node[anchor=north west,text=milcGradeText] at ([xshift=2.0cm,yshift=-1.55cm]current page.north west)
    {\sffamily\bfseries\fontsize{12}{14}\selectfont GRADO};
  \node[anchor=north west,text=milcGradeText,opacity=.85] at ([xshift=2.0cm,yshift=-2.35cm]current page.north west)
    {\sffamily\bfseries\fontsize{8.5}{10}\selectfont SERIE GUÍAS · TECNOLOGÍA E INFORMÁTICA};
  \node[anchor=north east,text=milcGradeText,opacity=.85,align=right] at ([xshift=-2.0cm,yshift=-1.55cm]current page.north east)
    {\sffamily\bfseries\fontsize{9}{12}\selectfont I.E. Sor María Juliana\\Cartago, Valle del Cauca};

  \node[anchor=west,text=milcGradeText] (gradonum) at ([xshift=1.85cm,yshift=-7.1cm]current page.north west)
    {\sffamily\bfseries\fontsize{112}{112}\selectfont 8};
  % El circulito de grado (°) solo aplica a las guias de grado («11°»); en
  % Bebras y Semillero el numero grande es un momento/modulo, no un grado.
  % Se posiciona relativo al nodo `gradonum`, no en coordenadas absolutas.
  \draw[milcGradeText,line width=1.6pt]
    ([xshift=2.0mm,yshift=-4.5mm]gradonum.north east) circle (.15cm);

  \node[anchor=west,text=milcGradeText,text width=11.4cm,align=left] at ([xshift=1.5cm,yshift=0.5cm]gradonum.east)
    {
     {\sffamily\bfseries\fontsize{9}{11}\selectfont\MakeUppercase{Período 2 · Lógica, algoritmos y computación física}}\\[3mm]
     {\sffamily\bfseries\fontsize{21}{25}\selectfont\setlength{\baselineskip}{27pt}Luces y sonido\\[4pt]con compás}\\[3.5mm]
     {\sffamily\bfseries\fontsize{15}{18}\selectfont Guía 15-8-TIC}
    };
\end{tikzpicture}

\vspace*{9.4cm}
\vfill

{\sffamily\bfseries\fontsize{8}{10}\selectfont\textcolor{milcGradeDark}{LO QUE TE LLEVAS HOY}}

\begin{tcolorbox}[enhanced,colback=white,colframe=milcNegro,arc=2mm,boxrule=1.6pt,
  left=5mm,right=5mm,top=4mm,bottom=4mm,before skip=3pt,after skip=12pt,
  fontupper=\RaggedRight\fontsize{11}{15.5}\selectfont]
Un programa en MakeCode con una animación de al menos ocho cuadros en la pantalla de 5×5, un tono de al menos cuatro notas sincronizado con los cuadros, y control con botones: A la inicia y B la detiene. Y una bitácora de tres versiones que diga qué cambiaste en cada una y por qué.
\end{tcolorbox}

\begin{tcolorbox}[enhanced,colback=milcGris,colframe=milcGris,arc=2mm,boxrule=0pt,
  left=5mm,right=5mm,top=3.5mm,bottom=3.5mm,after skip=12pt]
{\sffamily\bfseries\fontsize{8}{10}\selectfont\textcolor{milcNegro!55}{ESTA GUÍA ES DE}}\\[3mm]
\begin{tabularx}{\linewidth}{X p{3.2cm} p{3.2cm}}
\linead{\linewidth} & \linead{3.0cm} & \linead{3.0cm}\\[-1mm]
{\fontsize{8.5}{10}\selectfont\textcolor{milcNegro!55}{Nombre completo}} &
{\fontsize{8.5}{10}\selectfont\textcolor{milcNegro!55}{Grupo}} &
{\fontsize{8.5}{10}\selectfont\textcolor{milcNegro!55}{Fecha}}\\
\end{tabularx}
\end{tcolorbox}

\vfill

\noindent\textcolor{milcLinea}{\rule{\linewidth}{1pt}}\\[2mm]
{\sffamily\bfseries\fontsize{9.5}{12}\selectfont Autor: Dr. Álvaro Cárdenas Orozco}\\[1mm]
{\sffamily\fontsize{8.5}{11}\selectfont\textcolor{milcNegro!55}{Metodología MILC · apertura ancestral · pensamiento computacional · triángulo Dussel–estoicismo–Floridi}}

\newpage

\section*{Apertura: Actuadores --- luces y sonido con compás}

\begin{softbox}{milcGradeDark}{milcGradeSoft}{Saber ancestral}
En los ríos del Pacífico sur, entre Buenaventura y Tumaco, la gente se ha movido durante siglos en canoa, remando. Y mientras rema, canta. Los \textbf{cantos de boga} son cantos de trabajo: se entonan al ritmo de los remos, en trayectos que pueden durar un día entero (Ministerio de Cultura, 2010). El canto no quita el esfuerzo. Le pone un \textbf{compás}: cada golpe de remo cae donde cae la voz, y el cuerpo aguanta más cuando el esfuerzo tiene ritmo. Una animación en el micro:bit funciona igual. Ocho cuadros sin pausa son un parpadeo que nadie ve; ocho cuadros con una nota entre cada uno son una secuencia que cuenta algo. La \textbf{cara de exclusión}: remar durante días es un trabajo extenuante y mal pagado; el canto lo aligera, no lo dignifica por sí solo (Lozano Mayo y Palacios, 2013). Y una precisión: los cantos de boga son de los ríos del Pacífico, no del río Cauca ni del norte del Valle. Hoy vas a darle compás a tus luces con pausas y notas, para que la secuencia cuente algo.\par\smallskip{\footnotesize\textcolor{milcNegro!70}{\textbf{Fuente.} Ministerio de Cultura. (2010). \emph{Plan Especial de Salvaguardia de las músicas de marimba y los cantos tradicionales del Pacífico sur de Colombia} (A. Vanín Romero, coord.), p.~15. · Lozano Mayo, L. A. y Palacios, L. E. (2013). Los bogas del Chocó. \emph{Bioetnia}, 10, 61--69.}}
\end{softbox}

\begin{softbox}{milcTurquesa}{milcGris}{Saber contemporáneo}
Un \textbf{actuador} es lo que ejecuta la decisión del programa en el mundo: una luz, un sonido, un motor. El micro:bit trae dos. La \textbf{pantalla de 5×5}, 25 luces que se encienden una a una o en figuras; con el bloque «mostrar LEDs» pintas cada cuadro. Y el \textbf{altavoz}, en la versión 2, que toca notas con «tocar tono»: do, re, mi, con una duración. La animación es una secuencia de cuadros. Sin \textbf{pausa} entre ellos, cambian tan rápido que solo se ve el último. Con «pausa (ms)» entre cuadro y cuadro, se ve el movimiento. Y si la pausa dura lo mismo que la nota, la luz y el sonido van juntos: eso es sincronizar. Un programa de actuadores se construye por versiones. Primero los cuadros con pausa. Después las notas. Al final los botones que inician y detienen.
\end{softbox}

\begin{softbox}{milcMostaza}{milcGris}{Pregunta-puente}
\textit{El boga canta al ritmo del remo y por eso aguanta el trayecto. Cuando pongas ocho cuadros seguidos en el micro:bit sin pausa, ¿qué va a ver quien lo mire? ¿Y qué cambia si cada cuadro dura lo que dura una nota?}
\end{softbox}

\begin{softbox}{milcVerde}{milcGris}{Saber y saber hacer}
Al terminar podrás: (1) \textbf{identificar} qué actuadores tiene el micro:bit y qué bloque usa cada uno; (2) \textbf{crear} una animación de ocho cuadros con pausas y cuatro notas sincronizadas; (3) \textbf{aplicar} los botones A y B para iniciar y detener la secuencia; (4) \textbf{evaluar} cada versión y decir qué cambió y por qué.
\end{softbox}



\small
\begin{tabularx}{\textwidth}{>{\bfseries}p{3.0cm}Y}
\rowcolor{milcGradePrimary!12}Autor & Dr. Álvaro Cárdenas Orozco\\
DBA & Pensamiento computacional aplicado a sistemas de control con sensores y actuadores (MEN, T\&I)\\
Referentes & Pensamiento computacional · MakeCode / micro:bit · Logos como razón universal\\
Producto final & Un programa en MakeCode con una animación de al menos ocho cuadros en la pantalla de 5×5, un tono de al menos cuatro notas sincronizado con los cuadros, y control con botones: A la inicia y B la detiene. Y una bitácora de tres versiones que diga qué cambiaste en cada una y por qué.\\
\end{tabularx}
\normalsize

\puente{En la sesión 4 el micro:bit sintió y avisó. Hoy se expresa: luces en secuencia, sonido con compás y botones que mandan. En la sesión 6 vas a unir las dos cosas con variables y umbrales.}

\milcestacion{milcGradeDark}{Ruta MILC: así vamos a aprender}
\begin{tabularx}{\textwidth}{>{\centering\arraybackslash}X>{\centering\arraybackslash}X>{\centering\arraybackslash}X>{\centering\arraybackslash}X}
\phasecard{1}{milcVerde}{milcGris}{Escucha\\[-1mm]\small Observo, escucho y pregunto.} &
\phasecard{2}{milcTurquesa}{milcGris}{Entender\\[-1mm]\small Sistematizo y ordeno.} &
\phasecard{3}{milcMagenta}{milcGris}{Hacer\\[-1mm]\small Construyo y aplico.} &
\phasecard{4}{milcVino}{milcGris}{Cierre\\[-1mm]\small Evalúo y reflexiono.}
\end{tabularx}


\puente{Antes de programar, vas a dibujar en papel cuatro cuadros de una animación. Es coreografía en cuadrícula, antes del código.}

\milcestacion{milcVerde}{Estación 1 de 3 · Escucha}

\begin{softbox}{milcVerde}{milcGris}{Escena para escuchar}
\textbf{Actividad 1 · IDENTIFICA --- La coreografía en papel} (15 min · individual). (1) Dibuja en el cuaderno cuatro cuadrículas de 5×5: son la pantalla del micro:bit. (2) Elige una animación simple: un corazón que late, una flecha que gira, una ola que sube, un cuadrado que crece. (3) Rellena con lápiz las luces encendidas de cada cuadro, en orden. (4) Numera los cuadros del 1 al 4. (5) Debajo de cada uno escribe una nota que le pondrías: grave, media o aguda.
\end{softbox}

{\sffamily\bfseries\fontsize{8}{10}\selectfont\textcolor{milcNegro!55}{MARCA CUANDO LO HAYAS HECHO}}
\milccheckrow{Tengo cuatro cuadrículas de 5×5 rellenas y numeradas.}
\milccheckrow{La animación cuenta algo, late, gira, sube o crece.}
\milccheckrow{Cada cuadro tiene una nota anotada.}

\begin{infoband}{milcVerde}


\textbf{Lo que va al cuaderno (Actividad 1 · IDENTIFICA).} \textbf{Título:} «Actividad 1 --- La coreografía en papel». \textbf{Formato:} cuatro cuadrículas de 5×5 numeradas, con las luces encendidas rellenas y una nota debajo de cada una. \textbf{Extensión:} media página. \textbf{Sabes que terminaste cuando} los cuatro cuadros están rellenos, se nota el movimiento al mirarlos en orden, y cada uno tiene su nota.
\end{infoband}

\puente{Ya tienes los cuadros en papel. Ahora aprendes los bloques de LEDs, de música y de pausa, y con tu pareja armas la versión 1 y la versión 2.}

\milcestacion{milcTurquesa}{Estación 2 de 3 · Entender}

\begin{softbox}{milcTurquesa}{milcGris}{Lo que vas a entender}
\textbf{Actividad 2 · CREA --- Versión 1 y versión 2: luces y luego sonido} (30 min · parejas). (1) Con tu pareja, en MakeCode, arrastren «al presionar el botón A» y adentro un bloque «mostrar LEDs» por cada cuadro de su papel; completen hasta ocho. (2) Entre cuadro y cuadro pongan «pausa (ms)» de 200. Prueben: esa es la versión 1. (3) Ajusten las pausas hasta que la animación se vea fluida y anoten el valor. (4) Después de cada cuadro, o cada dos, agreguen «tocar tono» con una nota y una duración igual a la pausa. Prueben: esa es la versión 2. (5) Anoten qué cambió de la versión 1 a la 2.
\end{softbox}

\milcpaso{1}{milcTurquesa}{\textbf{Los bloques de luz.} «Mostrar LEDs» pinta un cuadro de 5×5 al hacer clic en las luces. «Mostrar ícono» trae figuras hechas. «Borrar pantalla» apaga todo. Una secuencia de «mostrar LEDs» con pausa entre cada uno es una animación.}
\milcpaso{2}{milcTurquesa}{\textbf{Los bloques de sonido.} «Tocar tono (nota) durante (tiempo)» toca una nota el tiempo que digas: do es 262, re 294, mi 330, fa 349, sol 392, la 440, si 494. Notas cortas hacen ritmo; largas se prolongan. Necesita el altavoz de la versión 2 del micro:bit o unos audífonos en los pines.}
\milcpaso{3}{milcTurquesa}{\textbf{La pausa es el compás.} Sin pausa, los cuadros cambian tan rápido que solo se ve el último. Con 200 milisegundos entre cuadros, se ve el movimiento. Si la nota dura lo mismo que la pausa, luz y sonido van juntos. Ese es el golpe del remo.}
\milcpaso{4}{milcTurquesa}{\textbf{Cómo se hace, paso a paso.} (1) Abre el papel. (2) Un «mostrar LEDs» por cuadro. (3) Pausa entre cada uno. (4) Prueba. (5) Una nota por cuadro, con la misma duración. (6) Prueba. (7) Anota qué cambió.}

\begin{drawbox}{milcTurquesa}{Los actuadores, en una mirada}
\textbf{Mostrar LEDs:} un cuadro de 5×5. \\[1mm] \textbf{Pausa (ms):} el compás entre cuadros. \\[1mm] \textbf{Tocar tono:} una nota con duración. \\[1mm] \textbf{Sincronizar:} nota y pausa del mismo largo. \\[1mm] \textbf{Botones:} A inicia, B detiene.
\end{drawbox}

\begin{softbox}{milcMostaza}{milcGris}{Errores comunes y su corrección}
(1) \textbf{Cuadros sin pausa}: solo se ve el último. (2) \textbf{Nota sin duración}: queda como pitido o suena eterna. (3) \textbf{Luces al azar}: ocho cuadros que no cuentan nada. (4) \textbf{Sin botón de parar}: la animación no se puede detener. (5) \textbf{Pantalla llena}: 25 luces encendidas a la vez borran la figura.
\end{softbox}

\begin{infoband}{milcTurquesa}


\textbf{Lo que va al cuaderno (Actividad 2 · CREA).} \textbf{Título:} «Actividad 2 --- Versión 1 y versión 2». \textbf{Formato:} la lista de los ocho cuadros con su pausa, la lista de las notas con su duración, y una línea sobre qué cambió de la versión 1 a la 2. \textbf{Extensión:} media página. \textbf{Sabes que terminaste cuando} la animación se ve fluida, al menos cuatro notas suenan con los cuadros, y anotaste el valor de pausa que funcionó.
\end{infoband}

\puente{Ya tienes luz y sonido juntos. Toca la versión 3, con botones que inician y detienen, y la bitácora de lo que cambió de una versión a otra.}

\milcestacion{milcMagenta}{Estación 3 de 3 · Hacer}

\begin{softbox}{milcMagenta}{milcGris}{Cómo vas a construir tu producto}
\textbf{Actividad 3 · EVALÚA --- Versión 3: los botones mandan} (25 min · individual). (1) Crea una variable «activa». En «al presionar el botón A» ponla en verdadero; en «al presionar el botón B», en falso. (2) Mueve la animación a «para siempre», dentro de un «si activa entonces». Prueba: A la inicia y B la detiene. (3) Escribe la bitácora de las tres versiones: qué tenía cada una, qué cambió y por qué. (4) Muéstrale la animación a tu pareja sin decir nada y pídele que diga qué cuenta. (5) Anota si acertó y, si no, qué cuadro cambiarías.
\end{softbox}

\milcpaso{1}{milcMagenta}{\textbf{Tu entrega tiene tres piezas.} (1) El programa final con ocho cuadros, cuatro notas y control por botones, compartido con enlace o archivo .hex. (2) La bitácora de las tres versiones. (3) Lo que dijo tu pareja que contaba la animación.}
\milcpaso{2}{milcMagenta}{\textbf{Por qué tres versiones.} La versión 1 ya funciona: son luces con pausa. La 2 le agrega sonido. La 3 le da control. Cada una es un producto, no un borrador. Trabajar así evita el programa gigante que nunca funciona.}
\milcpaso{3}{milcMagenta}{\textbf{La variable que manda.} «Activa» es un interruptor: verdadero o falso. A lo enciende, B lo apaga, y el «si activa» dentro de «para siempre» decide si la animación corre. Sin ese «si», los botones no mandan nada.}
\milcpaso{4}{milcMagenta}{\textbf{La prueba del que mira.} Tu animación cuenta algo solo si otro lo ve sin que se lo digas. Si tu pareja dice «un corazón que late», funciona. Si dice «unas luces», cambia el cuadro que confunde.}

\begin{drawbox}{milcMagenta}{Antes de entregar}
\checkbox Ocho cuadros distintos con pausa entre ellos. \\[1mm] \checkbox Al menos cuatro notas con la misma duración que la pausa. \\[1mm] \checkbox Botón A inicia y botón B detiene, con la variable «activa». \\[1mm] \checkbox Bitácora de tres versiones con qué cambió y por qué. \\[1mm] \checkbox Mi pareja dijo qué cuenta la animación sin que se lo explicara. \\[1mm] \checkbox Programa compartido con enlace o archivo .hex.
\end{drawbox}

\begin{softbox}{milcMagenta}{milcGris}{Plantilla de trabajo}
\textbf{Versión 1.} \textit{Ocho cuadros con pausa de … ms.} \\[1mm] \textbf{Versión 2.} \textit{Cuatro notas con la misma duración.} \\[1mm] \textbf{Versión 3.} \textit{Variable activa, A enciende, B apaga.} \\[1mm] \textbf{Prueba.} \textit{Lo que dijo mi pareja que cuenta.} \\[1mm] \textbf{Cambio.} \textit{El cuadro que ajusté y por qué.}
\end{softbox}

\begin{infoband}{milcMagenta}


\textbf{Lo que va al cuaderno (Actividad 3 · EVALÚA).} \textbf{Título:} «Actividad 3 --- Versión 3, los botones mandan». \textbf{Formato:} tabla de 3 filas y 3 columnas (versión / qué tenía / qué cambió y por qué), más lo que dijo tu pareja que contaba la animación. \textbf{Extensión:} media página. \textbf{Sabes que terminaste cuando} A inicia y B detiene, las tres filas están llenas y anotaste lo que vio tu pareja.
\end{infoband}

\puente{Lo que entregas hoy: el programa con ocho cuadros, cuatro notas y control por botones, y la bitácora de tres versiones. La prueba de calidad: tu pareja mira la animación sin que le expliques nada y dice qué cuenta.}

\milcestacion{milcMostaza}{Producto de la sesión}
\textbf{Animación de ocho cuadros + cuatro notas sincronizadas + control por botones + bitácora de tres versiones}.
\begin{softbox}{milcMostaza}{milcGris}{Criterios de calidad}
(1) La animación tiene ocho cuadros distintos con pausa y cuenta algo reconocible. (2) Al menos cuatro notas suenan con los cuadros, con la misma duración que la pausa. (3) El botón A inicia y el botón B detiene, con una variable. (4) La bitácora dice qué tenía cada versión y qué cambió. (5) Tu pareja dijo qué cuenta la animación sin que se lo explicaras. (6) El programa está compartido con enlace o archivo .hex.
\end{softbox}

\puente{Antes de cerrar, mira lo que hiciste desde cinco lados. Tres versiones que funcionan valen más que una perfecta que nunca terminaste.}

\milcestacion{milcVino}{Evaluación liberadora · cinco dimensiones}

\begin{softbox}{milcVino}{milcGris}{Sentido de la evaluación}
Evaluar no es solo revisar una nota. Es reconocer qué aprendí, qué cambió en mí, qué emociones manejé, cómo me relaciono con la ciudad y la comunidad, y qué le dejaré a quien viene detrás.
\end{softbox}

\textbf{Mi reflexión liberadora en cinco dimensiones.} Completa cada frase iniciada:

\small
\begin{tabularx}{\textwidth}{>{\bfseries\RaggedRight\arraybackslash}p{3.4cm}Y>{\RaggedRight\arraybackslash}p{4.6cm}}
\rowcolor{milcVino}\color{white}Dimensión & \color{white}\bfseries Mi respuesta (frame starter) & \color{white}\bfseries Algo que descubrí\\
1. Personal & Lo que más cambió en mí fue \linead{6.5cm} & \linead{4.2cm}\\
2. Emocional & La emoción más fuerte fue \linead{2.5cm} y la regulé \linead{3cm} & \linead{4.2cm}\\
3. Ciudadana & En democracia esto importa porque \linead{6cm} & \linead{4.2cm}\\
4. Local (Cartago) & En mi barrio sería útil para \linead{6.5cm} & \linead{4.2cm}\\
5. Intergeneracional & A mi abuela/abuelo le diría \linead{6.5cm} & \linead{4.2cm}\\
\end{tabularx}
\normalsize

\begin{infoband}{milcVino}
\textbf{Para la próxima sustentación.} Vuelve a esta tabla en seis meses. Verás que las respuestas cambian — no porque lo aprendido cambie, sino porque tú cambias. La evaluación liberadora es una conversación contigo a través del tiempo.
\end{infoband}


\puente{Para terminar, tres ideas sobre el ritmo y las versiones, contadas con palabras de esta guía: Dussel sobre la tecnología con rostro humano, Séneca sobre la vida que se nos va mientras la aplazamos, Floridi sobre la atención como bien escaso.}

\milcestacion{milcGradeDark}{Tres ideas para cerrar · Dussel, un estoico y Floridi}

\begin{softbox}{milcVino}{milcGris}{Enrique Dussel · lente del nosotros}
\textit{No hay liberación sin una tecnología con rostro humano, diseñada desde la historia de la gente que la usa.}

{\footnotesize\itshape\color{milcNegro!70}--- idea tomada de Enrique Dussel, contada con palabras de esta guía. No es una cita textual.}

Un canto de boga lo hicieron personas que remaban para vivir. Una animación que cuenta algo, un latido, una ola, es tecnología con rostro; ocho luces al azar, no. Lo que le da rostro es que alguien la entienda.

\textbf{¿Mi animación la entiende alguien que no soy yo?}
\end{softbox}
\linead{\linewidth}\\[1mm]
\linead{\linewidth}

\begin{softbox}{milcMostaza}{milcGris}{Estoicismo · Séneca · cuidado interior}
\textit{Mientras la aplazamos, la vida se va. Solo el tiempo es nuestro, y se escapa.}

{\footnotesize\itshape\color{milcNegro!70}--- idea tomada de Séneca, contada con palabras de esta guía. No es una cita textual.}

La versión 1 funciona hoy. Esperar a tener la versión perfecta es aplazar. Tres versiones que corren valen más que una idea enorme que nunca se probó.

\textbf{¿Qué versión de algo mío nunca entregué por esperar a que fuera perfecta?}
\end{softbox}
\linead{\linewidth}\\[1mm]
\linead{\linewidth}

\begin{softbox}{milcTurquesa}{milcGris}{Luciano Floridi y la Onlife Initiative · lente de la infoesfera}
\textit{La atención de las personas es un bien finito, precioso y escaso.}

{\footnotesize\itshape\color{milcNegro!70}--- idea tomada de Luciano Floridi y la Onlife Initiative, contada con palabras de esta guía. No es una cita textual.}

Ocho cuadros con compás cuidan la atención de quien mira; luces sin pausa la gastan. Diseñar una pantalla es decidir qué merece la atención de otro.

\textbf{¿Qué de mi animación pide atención sin devolver nada?}
\end{softbox}
\linead{\linewidth}\\[1mm]
\linead{\linewidth}

\begin{infoband}{milcNegro}
\textbf{Nota para el docente.} Las tres ideas están contadas con palabras de esta guía a partir del pensamiento de cada autor; no son citas textuales. De 9.º en adelante se trabajan con la obra y su referencia.
\end{infoband}

\milcestacion{milcVino}{Mi autoevaluación · marca una casilla por fila}

{\small
\setlength{\extrarowheight}{2.2pt}
\begin{tabularx}{\textwidth}{>{\RaggedRight\arraybackslash\bfseries}p{3.0cm}YYY}
\rowcolor{milcVino}\color{white}\bfseries Criterio & \color{white}\bfseries Lo logré & \color{white}\bfseries En proceso & \color{white}\bfseries Necesito apoyo\\[.6mm]
Comprensión del tema & \milcopcion{Explico las ideas con mis palabras.} & \milcopcion{Reconozco las ideas, falta precisión.} & \milcopcion{Necesito repasar lo básico.}\\[.8mm]
\rowcolor{milcGris}Pensamiento computacional & \milcopcion{Usé los pilares al armar mi producto.} & \milcopcion{Usé algunos pilares.} & \milcopcion{Necesito releer la estación 2.}\\[.8mm]
Aplicación y producto & \milcopcion{Mi producto cumple los criterios.} & \milcopcion{Está armado, con ajustes pendientes.} & \milcopcion{Aún está en construcción.}\\[.8mm]
\rowcolor{milcGris}Reflexión 5 dimensiones & \milcopcion{Respondí cada una con honestidad.} & \milcopcion{Respondí algunas con detalle.} & \milcopcion{Mis respuestas son muy generales.}\\[.8mm]
Triángulo de pensamiento & \milcopcion{Conecté las 3 voces con mi trabajo.} & \milcopcion{Conecté al menos una.} & \milcopcion{No las relaciono con mi proceso.}\\[.8mm]
\end{tabularx}}

\vspace{3mm}
\textbf{Hoy entendí que:}\\[1mm]
\linead{\linewidth}\\[1mm]
\linead{\linewidth}

\textbf{Mi compromiso para la próxima sesión es:}\\[1mm]
\linead{\linewidth}\\[1mm]
\linead{\linewidth}

\begin{softbox}{milcMostaza}{milcGris}{Cita de despedida · Marco Aurelio}
\textit{``El obstáculo en el camino se vuelve el camino. Lo que estorba al camino se vuelve camino.''}\\[1mm]
Cada error de hoy, bien observado, se convierte en pieza del camino que pisarás mañana.
\end{softbox}

\small
\begin{tabularx}{\textwidth}{>{\bfseries}p{3.6cm}Y}
\rowcolor{milcGradeDark!12}Nombre completo: & \linead{10cm}\\
Fecha: & \linead{4cm} \quad Curso/grupo: \linead{4cm}\\
Mi firma: & \linead{9cm}\\
\end{tabularx}
\normalsize

\begin{infoband}{milcGradeDark}
\textbf{Para la próxima sesión.} Llega con tu coreografía y la bitácora de tres versiones. En la sesión 6 vas a usar variables y umbrales para que los sensores no solo disparen acciones, sino que respondan a niveles. La variable «activa» de hoy es la primera de muchas.
\end{infoband}

\end{document}
